\documentclass[12pt, a4paper]{article}

\usepackage[utf8]{inputenc}
\usepackage[T1]{fontenc}
\usepackage[spanish]{babel}
\usepackage{amsmath}
\usepackage{amssymb}
\usepackage{geometry}
\usepackage{booktabs}
\usepackage{hyperref}
\usepackage{float}

% Para los bloques CML-DEVS usamos alltt, que permite
% mezclar texto normal con comandos LaTeX libremente.
\usepackage{alltt}
\usepackage{xcolor}

\geometry{
    top=2.5cm,
    bottom=2.5cm,
    left=3cm,
    right=3cm
}

% Comandos de conveniencia para CML-DEVS
\newcommand{\kw}[1]{\textbf{#1}}          % palabras clave en negrita
\newcommand{\cmlR}{$\mathbb{R}$}
\newcommand{\cmlTime}{\textbf{Time}}
\newcommand{\cmlsigma}{$\sigma$}
\newcommand{\dint}{$\delta_{\mathit{int}}$}
\newcommand{\dext}{$\delta_{\mathit{ext}}$}
\newcommand{\clambda}{$\lambda$}
\newcommand{\cmlta}{\textbf{ta}}
\newcommand{\cmlcomment}[1]{\textit{\textcolor{gray}{#1}}}

% Entorno para bloques CML-DEVS (imita el estilo de la tesis Hollmann)
\newenvironment{cmldevs}{%
  \begin{alltt}\small%
}{%
  \end{alltt}%
}

% ----------------------------------------------------------------

\title{
    \large Simulación de Sistemas \\[0.5em]
    \LARGE \textbf{Proyecto de Simulación} \\[0.3em]
    \large Sistema de Tiempo Real: Controlador Simplificado de Bomba de Infusión \\[1em]
    \normalsize Especificación DEVS de los Modelos Atómicos
}

\author{
    Jhonatan Calle Galeano \and Hernan Jara
}

\date{}

% ================================================================
\begin{document}
\maketitle
\thispagestyle{empty}

% ----------------------------------------------------------------
\section{Generador de Órdenes Médicas}

El modelo \textbf{GeneradorOrdenes} es un modelo atómico autónomo (sin entradas)
que genera periódicamente eventos \texttt{ordenMedica}, indicando el caudal objetivo
que debe administrar la bomba. Las órdenes pueden generarse de manera determinística
o estocástica según los escenarios definidos.

\begin{figure}[H]
\begin{cmldevs}
\kw{atomic} GeneradorOrdenes \kw{is} <Y, S, \dint{}, \clambda{}, \cmlta{}> \kw{where}
\kw{S is}
    caudalProximo : \cmlR{};    \cmlcomment{-- caudal objetivo (0 a 200 ml/h)}
    interarribo   : \cmlTime{};  \cmlcomment{-- tiempo hasta la siguiente orden}
    \cmlsigma{}             : \cmlTime{};  \cmlcomment{-- variable de avance de tiempo}
\kw{end S}
\kw{Y is}
    ordenMedica : \cmlR{};      \cmlcomment{-- puerto de salida: valor del caudal}
\kw{end Y}
\dint{}(caudalProximo, interarribo, \cmlsigma{}) \kw{is}
    \cmlcomment{-- A definir en la implementacion: puede ser deterministica}
    \cmlcomment{-- (valores fijos) o estocastica (usando distribuciones)}
    caudalProximo = next\_caudal\_logic();
    interarribo   = next\_time\_logic();
    \cmlsigma{}   = interarribo;
\kw{end} \dint{}
\clambda{}(caudalProximo, interarribo, \cmlsigma{}) \kw{is}
    \cmlcomment{-- Emite el caudal objetivo por el puerto de salida}
    ordenMedica = caudalProximo;
\kw{end} \clambda{}
\cmlta{}(caudalProximo, interarribo, \cmlsigma{}) \kw{is}
    \cmlsigma{};
\kw{end} \cmlta{}
\kw{end atomic}
\end{cmldevs}
\caption{Especificación CML-DEVS: Generador de Órdenes Médicas}
\label{fig:generador-ordenes}
\end{figure}


\section{Sensor de Flujo}

El modelo \textbf{SensorFlujo} actúa como puente entre el actuador y el controlador,
reportando periódicamente el caudal real cada 1 segundo exacto.

\begin{figure}[H]
\begin{cmldevs}
\kw{atomic} SensorFlujo \kw{is} <Y, S, \dint{}, \clambda{}, \cmlta{}> \kw{where}
\kw{S is}
    caudalMedido : \cmlR{};  \cmlcomment{-- valor de caudal generado internamente}
    \cmlsigma{}          : \cmlTime{}; \cmlcomment{-- controla el periodo de muestreo (1 s)}
\kw{end S}
\kw{Y is}
    sensorFlujo  : \cmlR{};  \cmlcomment{-- medicion enviada al Controlador}
\kw{end Y}
\dint{}(caudalMedido, \cmlsigma{}) \kw{is}
    \cmlcomment{-- Genera el siguiente valor de caudal y reinicia el temporizador}
    caudalMedido = next\_caudal\_logic();
    \cmlsigma{}          = 1.0;
\kw{end} \dint{}
\clambda{}(caudalMedido, \cmlsigma{}) \kw{is}
    \cmlcomment{-- Emite el valor medido antes de la transicion interna}
    sensorFlujo = caudalMedido;
\kw{end} \clambda{}
\cmlta{}(caudalMedido, \cmlsigma{}) \kw{is}
    \cmlsigma{};
\kw{end} \cmlta{}
\kw{end atomic}
\end{cmldevs}
\caption{Especificación CML-DEVS: Sensor de Flujo}
\label{fig:sensor-flujo}
\end{figure}

\subsection*{Análisis del modelo}

\begin{enumerate}
    \item \textbf{Frecuencia de muestreo} ($\mathit{ta}$ y $\delta_{\mathit{int}}$): El requisito establece que el sensor debe informar el caudal real cada 1 segundo. Por ello, $\sigma$ se inicializa y resetea en \texttt{1.0} en cada transición interna.

    \item \textbf{Puertos} ($Y$): Se define el puerto de tipo $\mathbb{R}$ para la  comunicación modular dentro del modelo acoplado.

    \item \textbf{Estado inicial} ($S$): Se asume \texttt{caudalMedido} $= 0.0$ y $\sigma = 1.0$, indicando que la primera medición ocurrirá exactamente un segundo después del arranque.
\end{enumerate}

\section{Actuador de la Bomba}

El modelo \textbf{ActuadorBomba} representa la interfaz física de la bomba: recibe
comandos lógicos del controlador y, tras una latencia de 0.5 segundos, actualiza
el caudal real suministrado al paciente.

\begin{figure}[H]
\begin{cmldevs}
\kw{atomic} ActuadorBomba \kw{is} <X, Y, S, \dint{}, \dext{}, \clambda{}, \cmlta{}> \kw{where}
\kw{S is}
    caudalFisico   : \cmlR{};       \cmlcomment{-- caudal que actualmente sale de la bomba}
    caudalObjetivo : \cmlR{};       \cmlcomment{-- caudal a alcanzar tras la latencia}
    \cmlsigma{}            : \cmlTime{};    \cmlcomment{-- controla la latencia de 0.5 segundos}
\kw{end S}
\kw{X is}
    ajustarCaudal : \cmlR{};        \cmlcomment{-- nuevo valor de dosis desde el controlador}
    detenerBomba  : \{signal\};     \cmlcomment{-- senal de parada de emergencia}
\kw{end X}
\kw{Y is}
    caudalActual : \cmlR{};         \cmlcomment{-- notifica el cambio fisico al Sensor}
\kw{end Y}
\dext{}((caudalFisico, caudalObjetivo, \cmlsigma{}), e, (puerto, valor)) \kw{is}
    \kw{defcases}
        caudalObjetivo = valor;   \kw{if} (puerto = ajustarCaudal)
        \cmlsigma{} = 0.5;
        caudalObjetivo = 0.0;     \kw{if} (puerto = detenerBomba)
        \cmlsigma{} = 0.5;
    \kw{end defcases}
\kw{end} \dext{}
\dint{}(caudalFisico, caudalObjetivo, \cmlsigma{}) \kw{is}
    \cmlcomment{-- Cumplida la latencia, el caudal fisico alcanza al objetivo}
    caudalFisico = caudalObjetivo;
    \cmlsigma{}          = infinity;
\kw{end} \dint{}
\clambda{}(caudalFisico, caudalObjetivo, \cmlsigma{}) \kw{is}
    \cmlcomment{-- Emite el nuevo caudal real antes de pasar al estado pasivo}
    caudalActual = caudalObjetivo;
\kw{end} \clambda{}
\cmlta{}(caudalFisico, caudalObjetivo, \cmlsigma{}) \kw{is}
    \cmlsigma{};
\kw{end} \cmlta{}
\kw{end atomic}
\end{cmldevs}
\caption{Especificación CML-DEVS: Actuador de la Bomba}
\label{fig:actuador-bomba}
\end{figure}

\subsection*{Análisis del modelo}

\begin{enumerate}
    \item \textbf{Manejo de la latencia} ($\delta_{\mathit{ext}}$ y $\mathit{ta}$): Al recibir
    un evento externo, el modelo guarda el nuevo valor en \texttt{caudalObjetivo} y
    establece $\sigma = 0.5$, obligando al sistema a esperar exactamente el tiempo
    de latencia especificado antes de aplicar el cambio.

    \item \textbf{Estado pasivo}: Cuando no hay órdenes pendientes, $\sigma = \infty$.
    El actuador mantiene el caudal actual indefinidamente hasta que el controlador
    envíe un nuevo evento.

    \item \textbf{Sincronización con el sensor}: La salida \texttt{caudalActual} se
    conecta en el modelo acoplado a la entrada del Sensor de Flujo, que reportará
    el cambio medio segundo después de que el controlador emitió la orden.

    \item \textbf{Tipos de datos}: Se utiliza $\mathbb{R}$ para los caudales
    (0 a 200 ml/h) y \texttt{\{signal\}} para el evento discreto de detención.
\end{enumerate}
\end{document}
